\documentclass[spanish,openany]{report}
% \usepackage{darkmode} \enabledarkmode
\usepackage{wrapfig}
\input{~/git/Clase/template/preamble.tex}
\input{~/git/Clase/template/macros.tex}
\input{~/git/Clase/template/letterfonts.tex}
\graphicspath{{/home/kuma/git/Clase/Fp_Gs_Robotica/images/}}
\setimagefolder{/home/kuma/git/Clase/Fp_Gs_Robotica/images/}

\date{\today}
% \hbadness=99999
\usepackage{microtype}
% \addbibresource{./bibliografia.bib}
\usepackage{titling}
\usepackage{afterpage}


% Header and footer
% \fancypagestyle{plain}{
%     \fancyhf{} %Clear Everything.
%     \fancyfoot[C]{Página \thepage\ de \pageref{LastPage} }%Page Number
%     \renewcommand{\headrulewidth}{1pt} %0pt for no rule, 2pt thicker etc...
%     \renewcommand{\footrulewidth}{1pt}
%     \fancyfoot[R]{\small{Daniel Carbajales Soto}}
%     \fancyhead[R]{\small{Cinta Transportadora con 3 Velocidades}}
%     \fancyhead[L]{\small{\today}}
% }
% \pagestyle{plain}

\begin{document}

%----------------------------------------------------------------------------------------
%	                            PORTADA
%----------------------------------------------------------------------------------------
% \makeportada{}{Daniel Carbajales Soto}{ 1° CFGS Automatización y Robótica Industrial}{}

%----------------------------------------------------------------------------------------
%                                    HEADER 
%----------------------------------------------------------------------------------------
% \header{}{Daniel Carbajales Soto}{1° CFGSARI}
%----------------------------------------------------------------------------------------
%                               TABLE OF CONTENTS
%----------------------------------------------------------------------------------------
% \maketoc
%----------------------------------------------------------------------------------------
%	                            TABLES
%----------------------------------------------------------------------------------------
% this is just a demo in case i need to make a table
%
%
%
% \begin{tabularx}{\linewidth}{|C{2cm}|Y|Y|C{2cm}|}
% \hline
% \makecell{Header 1} & \makecell{Header 2 \\ with linebreak} & \makecell{Header 3} & \makecell{Header 4} \\
% \hline
% Row 1 & This is a long cell text that automatically wraps in the Y column & Another long cell & Short \\
% \hline
% Row 2 & \multirow{2}{*}{Multirow text} & More text & Short \\
% \cline{1-1} \cline{3-4}
% Row 3 & & Extra text & Short \\
% \hline
% \end{tabularx}


%----------------------------------------------------------------------------------------
%	                            DOCUMENT
%----------------------------------------------------------------------------------------


\section*{Practica 2}
\subsection*{Objetivo:}
Estudiar el efecto de una inductancia y una resistencia conectadas en serie. 
Determinar la reactancia inductiva y determinar la impedancia. Ver el coseno de $\varphi $ con un Fasímetro.

\subsection*{Material:}
\begin{minipage}[t]{0.48\textwidth}
\begin{itemize}[label={\textbullet}, leftmargin=*, nosep]
\item 1 lámpara de incandescencia y 1 reactancia.
\item 1 Voltímetro.
\item 1 Vatímetro.
\item 1 Polímetro digital
\end{itemize}
\end{minipage}%
\hfill
\begin{minipage}[t]{0.48\textwidth}
\begin{itemize}[label={\textbullet}, leftmargin=*, nosep]
\item 1 Panel de pruebas.
\item 1 Amperímetro.
\item 1 Fasímetro.
\end{itemize}
\end{minipage}

\subsection*{Fundamento Teórico:}
En los circuitos de continua, las resistencias limitan la corriente. De igual manera,
las resistencias se oponen al paso de la misma en los circuitos do alterna. Existen además 
otros componentes, denominados reactivos, que se oponen al paso de la corriente en los 
circuitos de alterna: las inductancias (bobinas de autoinducción) y los condensadores. 
Los efectos reactivos de estos componentes no pueden ser determinados por medio del óhmetro.
La propiedad que tiene una bobina de oponerse a cualquier variación de la corriente que la 
atraviesa, es denominada autoinducción (L) y se mide en Henrio (H).
Para determinar la corriente alterna que circula por una bobina habrá pues que hallar la 
autoinducción de la misma, además de su resistencia en corriente continua. 
La oposición que presenta una inductancia a la corriente alterna se denomina reactancia inductiva, 
su valor se expresa en ohmios y viene dada por la fórmula: 
\[
	X_{L} = 2\pi fL
\]




























%%%%%%%%%%%%%%%%%%%%%%%%%%%%%%%%%%%%%%%%%%%%%%%%%%%%%%%%%%
%%%%%%%%%%%%%%%%%%%% bibliography %%%%%%%%%%%%%%%%%%%%%%%%
% \nocite{*}                                 %%%%%%%%%%%%%
% \printbibliography                         %%%%%%%%%%%%%
%%%%%%%%%%%%%%%%%%%%%%%%%%%%%%%%%%%%%%%%%%%%%%%%%%%%%%%%%%
%%%%%% Last page in case there is no bibliography %%%%%%%%
% \label{Last Page}                           %%%%%%%%%%%%%
%%%%%%%%%%%%%%%%%%%%%%%%%%%%%%%%%%%%%%%%%%%%%%%%%%%%%%%%%%

\end{document}
